\section{25维：接触层的协同变形}\label{sec:motion25}

固定构型不能直接补点，并不排除移动其中一部分后获得增益。
下面的构造从公开的197579点Leech提升出发\cite{kissingnumbers}，
移动552个保留赤道点，加入一个新点，再修复两个提升点。
协同移动保持块内的全部距离，与未动赤道的相容性则由凸性保证。

\begin{lemma}[沿弦的协同移动]\label{lem:chord-motion}
设 $n$ 为单位向量，$h>0$，且 $r_i\perp n$、
$\|r_i\|^2=\rho-h^2$。令 $u_i=r_i-hn$、$w_i=r_i+hn$。
对满足 $c^2+d^2=h^2$ 的 $c,d$，定义 $W_i=(r_i+cn,d)$。
则 $\|W_i\|^2=\rho$，且 $W_i\cdot W_j=w_i\cdot w_j$。
若 $z\cdot u_i\le\rho/2$ 且 $z\cdot w_i\le\rho/2$，
则 $(z,0)\cdot W_i\le\rho/2$。
\end{lemma}
\begin{proof}
正交关系给出范数与内积恒等式。由于 $|c|\le h$，水平分量
$\frac{h-c}{2h}u_i+\frac{h+c}{2h}w_i$ 是两个端点的凸组合，
故与 $z$ 的内积保持相同上界。
\end{proof}

对有限非空移动层，固定点 $(z,s)$ 给出半平面约束
\[
c\langle n,z\rangle+ds\le\rho/2-\max_i\langle r_i,z\rangle.
\]
因此，$m$ 项边界约束的可行性只需检查至多 $2m+1$ 个参数：
圆与各边界直线的交点，再加 $(h,0)$。在这些交点之间的每段开弧上，
所有不等式的真假不变；移到端点不会破坏已满足的闭半平面约束。
同理，对可删除约束赋予非负代价后，某个候选点达到最小删除代价。
有理边界数据与有理 $h^2$ 只需二次代数数计算。
轴向新增点 $Q=(\alpha n,\beta)$、$\alpha^2+\beta^2=\rho$ 与全层接触，
则另要求 $\alpha c+\beta d=\rho/2$；该直线与圆有交点当且仅当
$h^2\ge\rho/4$。候选点与未动点的相容性仍须单独检查。

\subsection{伙伴层与精确移动}
采用平方范数4、不同点内积至多2的归一化。
设 $L$ 为196560点Leech极小壳，$D$ 为原提升删除的1016个方向，
$x_i$ 为原1016个头。原构型为
\[
\{(z,0):z\in L\setminus D\}\ \cup\
\{(x_i,\pm1):0\le i<1016\}\ \cup\ \{(P,0),(0,2),(0,-2)\}.
\]
原头数据、删除方向及其精确坐标保存于证书。
取 $n=v/\sqrt6=-P/2$，其中
\[
\sqrt8\,v=(-5,-1,-1,-1,-1,1,-1,1,-1,-1,-1,1,
1,-1,-1,-1,1,1,-1,1,-1,-1,1,-1).
\]
集合 $U=\{u\in L:u\cdot v=-3\}$ 有552个点，全部属于 $D$。
每个 $u$ 的伙伴 $w=u+v$ 是未删除的极小向量，且 $u\cdot w=1$。
令
\[
r=u+\sqrt{3/2}\,n=w-\sqrt{3/2}\,n,\qquad
a=\frac{\sqrt3}{2}+\frac{\sqrt2}{4},\qquad b=\frac{\sqrt6-2}{4}.
\]
于是 $r\perp n$、$r^2=5/2$。将每个 $(w,0)$ 换为
$W(r)=(r+an,b)$，并加入 $Q=(\sqrt3n,-1)$。

\paragraph{解析点对检查。}
沿弦引理给出移动层的范数、内部内积及其与未动赤道的相容性。
恒等式
\[
a^2+b^2=3/2,\qquad \sqrt3a-b=2,\qquad -a/\sqrt2+b=-3/4
\]
给出 $Q\cdot W(r)=2$。
Leech格的整性及最小范数给出 $v\cdot z\in\{-3,-2,-1,0,1,2,3\}$；
负极端层已删除，正极端层恰为移动层。因此 $Q$ 与未动赤道的
内积至多 $\sqrt2$。

与 $(P,0)$、两极的内积分别为
$-2a,\pm2b$（对于 $W$）及 $-2\sqrt3,-2,2$（对于 $Q$）。
完整原块的头为 $r-n/\sqrt2$。
同一块编号上的上、下提升点与 $W$ 的内积分别为
$7/4$ 和 $11/4-\sqrt6/2$，均小于2。
不同块编号有 $r_i\cdot r_j\le1/2$，相应内积更小。
\subsection{两个局部修复与有限检验}
替换上侧点 $(x_{984},1)$ 和下侧点 $(x_{1008},-1)$，
保留它们另一侧的伙伴。两条替换向量为
\[
R_+=2k/\sqrt{1007176},\qquad R_-=2\ell/\sqrt{10209},
\]
其中整数行 $k,\ell$ 由\texttt{constructions/d25/repair-points.json}指定，
平方范数分别为1007176和10209，故 $R_\pm^2=4$。

固定 $Q$、保持完整沿弦移动层时，至少需改动两个原帽点；所给方案达到这个最小数目。
首先 $Q\cdot(x_{1008},-1)>231/100>2$，所以该下侧点必须改变。
若只改它，所有第一块的成对帽点仍在，迫使
$-c/\sqrt2+|d|\le-1/2$，结合圆方程即得 $c>0$、$|d|\le1/3$。
在该区间内 $\sqrt3\sqrt{3/2-d^2}-d$ 严格递减，故与 $Q$ 相容要求 $d\ge b$。
令 $\nu=\langle n,x_{984}\rangle$、$M=\max_r\langle r,x_{984}\rangle$。
上侧984号点与移动层的最大内积为
\[
g(d)=M+\nu\sqrt{3/2-d^2}+d,\qquad
g'(d)\ge1-\sqrt6/5>0\quad(b\le d\le1/3).
\]
精确头数据的外向舍入计算给出 $g(b)>201/100>2$，故该点也必须改变。
运动障碍检查器复核这两项严格不等式（见附录~\ref{app:data}）。

余下有限检查包括552个移动点与2030个未替换提升点、
$Q$ 与这些提升点，以及两个修复点与最终其余全部点的内积。
192位Arb区间计算严格确认四组上界分别小于
\[
1.999372,\qquad1.965926,\qquad1.994722,\qquad1.987522,
\]
均低于2。原构型的头--壳及头--头检查使用整数和有理数。
移动产生的等号情形由前述引理处理，剩余未动点对沿用原构型的相容性。

\begin{theorem}
上述数据确定197580个不同的25维接吻点，故 $K(25)\ge197580$。
\end{theorem}
\begin{proof}
上述解析与有限检查覆盖所有点对。552个移动点及两个修复点均是等数替换，
只有 $Q$ 增加点数，故
\[
197579-552-2+552+2+1=197580.
\]
若两条不同索引的向量相等，其内积将为4，与已验证的上界2矛盾。
\end{proof}

\texttt{constructions/d25/}给出精确原始数据、修复向量、完整独立说明与程序。
\texttt{dimension25}复现任务同时检查原构型、扩展和两项破坏证书测试。
